\section{Key Findings and Clinical Takeaways}

This section synthesizes the evidence from 13 primary studies and 5 exploratory feasibility studies on wearable seizure detection and forecasting, identifying patterns that inform clinical practice and future research directions.

\subsection{Executive Summary}

\begin{itemize}
  \item \textbf{Current State:} Wearable seizure detection achieves high sensitivity (71--100\%) for convulsive seizures using accelerometry, but false alarm rates remain clinically unacceptable in most studies.

  \item \textbf{Forecasting Challenge:} Seizure forecasting remains immature, with only 43.5--83\% of patients achieving better-than-chance performance and significant heterogeneity in individual responses.

  \item \textbf{Best Performing Approaches:} Multimodal sensor fusion (ACC + cardiac) and ensemble deep learning methods show the most promise for clinical translation.

  \item \textbf{Validation Gap:} Only 3 of 13 studies include prospective validation, with most relying on retrospective analyses that may overestimate real-world performance.
\end{itemize}

\subsection{Detection Studies: Performance Landscape}

\subsubsection{High-Sensitivity Detectors}
Table~\ref{tab:detection-performance} summarizes the best-performing detection systems across different seizure types and settings.

\begin{table}[htbp]
\centering
\small
\caption{Highest Performing Detection Studies by Metric}
\label{tab:detection-performance}
\begin{tabular}{llccccc}
\toprule
\textbf{Study} & \textbf{Seizure Type} & \textbf{Sensitivity} & \textbf{FAR} & \textbf{Latency} & \textbf{Validation} & \textbf{Ready} \\
\midrule
Fine 2025 & Tonic & 100\% & 0.023/h & 14~s & Phase 1 & $\circ$ \\
Dong 2026 & Nocturnal severe & 71.6\% & 0.165/h & $<$30~s & Prospective & $\bullet$ \\
Spahr 2025 & Convulsive & 96\% & $<$1/8d & 26~s & Phase 2 & $\bullet$ \\
Reintjes 2025 & All types & 98.16\% & 1.91/h & -- & Retro & $\circ$ \\
Borujeny 2013 & Motor & 100\% & 0 & -- & 3-pt & $\circ$ \\
\bottomrule
\end{tabular}
\end{table}

\noindent \textbf{Key Finding:} Only two studies (Dong, Fine) achieve false alarm rates approaching clinical targets ($<$0.5/h) while maintaining reasonable sensitivity. The Phase 2 Spahr study demonstrates that multi-center validation is feasible but still shows FAR above clinical targets.

\subsubsection{Clinical Translation Barriers}
\begin{enumerate}
  \item \textbf{False Alarm Crisis:} Reintjes et al. report FAR of 1.91--39.75/h, 20--200$\times$ higher than clinically acceptable. This confirms that ECG-only approaches without motion artifact filtering are not viable for standalone deployment.

  \item \textbf{Seizure Type Limitation:} High sensitivity ($>$90\%) is consistently achieved only for convulsive/motor seizures. Non-motor seizures remain largely undetectable by current wearable approaches.

  \item \textbf{Offline Processing:} Most algorithms (Fine, Wang, Singh Rathore) run on offline servers rather than embedded devices, limiting real-time applicability.
\end{enumerate}

\subsection{Forecasting Studies: The Prediction Challenge}

\subsubsection{Performance Spectrum}
Forecasting performance varies dramatically across studies and patients (Table~\ref{tab:forecasting-performance}).

\begin{table}[htbp]
\centering
\small
\caption{Forecasting Study Performance Comparison}
\label{tab:forecasting-performance}
\begin{tabular}{lccccc}
\toprule
\textbf{Study} & \textbf{Metric} & \textbf{Value} & \textbf{TiW} & \textbf{Patient Success} & \textbf{Lead Time} \\
\midrule
Meisel 2020 & IoC & 14.1\% & 43.7\% & 43.5\% (30/69) & -- \\
Nasseri 2021 & AUC-ROC & 0.75 & 0.9--7.2~h/d & 83\% (5/6) & 33~s mean \\
Stirling 2021 & AUC & 0.74 & -- & 100\% (hourly) & Cyclic HR \\
Vieluf 2025 & Sensitivity & 82\% & -- & -- & Daily \\
Ode 2023 & Sensitivity & 74\% & -- & 44\% (29/66) & -- \\
\bottomrule
\end{tabular}
\end{table}

\noindent \textbf{Critical Finding:} Only Meisel et al. used true LOSO cross-validation, revealing that patient-specific tuning inflates performance by 30--40 percentage points. Their 43.5\% patient success rate suggests that fewer than half of patients may benefit from current forecasting approaches.

\subsubsection{The Responder Problem}
\begin{itemize}
  \item Meisel (2020): 43.5\% significant forecasters
  \item Ode (2023): 44\% patient success (29/66 with 100\% sensitivity)
  \item Nasseri (2021): 83\% success (5/6 above chance)
  \item Stirling (2021): 100\% hourly success but highly variable AUC
\end{itemize}

\noindent \textbf{Interpretation:} Patient heterogeneity is the fundamental challenge in seizure forecasting. Universal models (Vieluf) may report aggregated performance that masks significant individual variation.

\subsection{Methodological Insights}

\subsubsection{Algorithm Evolution Trajectory}
\begin{table}[htbp]
\centering
\small
\caption{Evolution of Algorithms in Seizure Detection/Forecasting (2013--2026)}
\label{tab:algorithm-evolution}
\begin{tabular}{llll}
\toprule
\textbf{Period} & \textbf{Algorithms} & \textbf{Representative Studies} & \textbf{Characteristics} \\
\midrule
2013--2015 & KNN, ANN, SVM & Borujeny 2013, Fujiwara 2016 & Handcrafted features, small samples \\
2016--2019 & HRV-based MLP, LSTM & Yamakawa 2020, Billeci 2018 & Patient-specific tuning emerges \\
2020--2022 & LSTM+1DConv, Ensembles & Meisel 2020, Stirling 2021 & LOSO validation, multifactor \\
2023--2026 & CNN+Attention, Matrix Profile, Self-Att AE & Ode 2023, Reintjes 2025, Spahr 2025 & Anomaly detection, \\
& & & attention mechanisms \\
\bottomrule
\end{tabular}
\end{table}

\subsection{Exploratory Studies: Feasibility Evidence}

\subsubsection{Cardiac Signal Viability}
The five exploratory studies collectively demonstrate that cardiac signals (ECG, HRV, PPG) contain detectable preictal changes:

\begin{itemize}
  \item \textbf{Pavei 2017:} 94.1\% sensitivity using SVM on HRV dynamics (n=12)
  \item \textbf{Fujiwara 2016:} 91\% sensitivity using multivariate statistical process control
  \item \textbf{Leal 2017:} Identified RRMean and BPMMean as most consistent ECG features (n=167)
  \item \textbf{Jain 2017:} Demonstrated EDA decomposition feasibility for seizure-related autonomic changes
  \item \textbf{Seth 2023:} Systematic review showing 56--100\% sensitivity range across 24 studies
\end{itemize}

\noindent \textbf{Synthesis:} Cardiac signals are viable for seizure detection but require sophisticated signal processing to overcome motion artifacts and individual variability.

\subsection{Clinical Recommendations}

\subsubsection{For Current Clinical Practice}
Based on current evidence, the following recommendations can be made:

\begin{enumerate}
  \item \textbf{Nocturnal Monitoring:} Dong et al.'s NightWatch system (71.6\% sensitivity, 0.165/h FAR) is the most validated option for home nocturnal monitoring.

  \item \textbf{Tonic Seizures:} Fine et al.'s algorithm (100\% sensitivity, 0.023/h FAR) shows promise for tonic seizure detection but requires Phase 2 validation.

  \item \textbf{EMU Settings:} Spahr et al.'s Phase 2 validated algorithm (96\% sensitivity, $<$1/8d FAR) is appropriate for EMU seizure counting.

  \item \textbf{Avoid:} Standalone ECG-based detection systems (FAR $>$1/h) without confirmation from accelerometry.
\end{enumerate}

\subsubsection{For Researchers and Developers}
\begin{enumerate}
  \item \textbf{Validation Standard:} LOSO cross-validation should be the minimum standard for reporting performance.

  \item \textbf{Metric Reporting:} All studies should report: sensitivity, specificity/PPV, FAR per hour, detection latency, and patient success rate.

  \item \textbf{Interpretability:} SHAP or similar explainable AI frameworks should be incorporated to enable clinical trust and regulatory approval.

  \item \textbf{Responder Identification:} Pre-study characterization of likely responders based on seizure type, HRV reactivity, and other biomarkers.
\end{enumerate}

\subsection{Limitations of Current Evidence}

\subsubsection{Study Design Limitations}
\begin{itemize}
  \item \textbf{Small Samples:} Median n=15--70 patients; most studies underpowered for subgroup analysis
  \item \textbf{Short Duration:} Median monitoring $<$1 week per patient; rare events (seizures) limit statistical power
  \item \textbf{Retrospective Bias:} 10/13 studies retrospective; high risk of data leakage and optimistic bias
  \item \textbf{Setting Bias:} Most studies in EMU; home performance likely worse
\end{itemize}

\subsubsection{Metric Heterogeneity}
\begin{itemize}
  \item Sensitivity vs. accuracy vs. AUC reported inconsistently
  \item FPR reported as per-hour, per-day, per-night, or not at all
  \item Patient success rate reported by only 6/13 studies
  \item No standardized benchmark datasets (except EPILEPSIAE)
\end{itemize}

\subsection{Future Directions}

\subsubsection{Research Priorities}
Based on identified gaps, the following research priorities emerge:

\begin{enumerate}
  \item \textbf{Standardized Metrics:} Develop consensus on minimum reporting standards (ILAE Phase-appropriate metrics)

  \item \textbf{Responder Prediction:} Identify baseline biomarkers that predict which patients will achieve usable forecasting performance

  \item \textbf{Edge Optimization:} More studies like Spahr (1.12~s inference) needed for real-time deployment

  \item \textbf{Personalization at Scale:} Federated learning approaches that enable patient-specific models without data centralization

  \item \textbf{Hybrid Systems:} Combined wearable + diary approaches (Vieluf) that leverage complementary information sources
\end{enumerate}

\subsubsection{Emerging Technologies}
\begin{itemize}
  \item \textbf{Matrix Profile Anomaly Detection:} Reintjes et al. show unsupervised approaches can achieve high sensitivity without seizure-specific training
  \item \textbf{Self-Attentive Autoencoders:} Ode et al. demonstrate AE approaches for RRI-based forecasting
  \item \textbf{Attention Mechanisms:} Dong et al. use attention to weight reliable signals and reduce motion artifact impact
  \item \textbf{Ensemble Deep Learning:} Spahr et al. show ensemble CNN approaches achieve Phase 2 validation
\end{itemize}

\subsection{Conclusion}

Wearable seizure detection and forecasting have advanced significantly from proof-of-concept studies to Phase 2 clinical validation. Detection of convulsive seizures is now reliable enough for limited clinical use, with two systems (Dong, Fine) achieving false alarm rates compatible with patient adoption. However, seizure forecasting remains an open challenge, with fewer than half of patients achieving better-than-chance performance in LOSO-validated studies.

The path forward requires: (1) rigorous LOSO validation as the standard, (2) standardized metric reporting enabling cross-study comparison, (3) interpretable models that earn clinician trust, and (4) recognition that patient heterogeneity may require personalized approaches rather than universal solutions. The integration of multimodal sensors, attention mechanisms, and explainable AI represents the current state-of-the-art, but clinical translation demands prospective, multi-center validation in real-world settings.
